\documentclass[reqno,11pt]{amsart}

\usepackage[margin=1in]{geometry}
\usepackage[T1]{fontenc}
\usepackage{lmodern}
\usepackage{microtype}
\usepackage{amsmath,amssymb,amsthm,mathtools}
\usepackage[dvipsnames]{xcolor}
\usepackage[most]{tcolorbox}
\usepackage{tikz}
\usepackage{enumitem}
\usepackage[hidelinks]{hyperref}
\hypersetup{
  pdftitle={Lecture 3, MATH 327: Permutation Groups},
  pdfauthor={Manish Patnaik},
  pdfsubject={MATH 327 course notes}
}

\definecolor{courseblue}{HTML}{1F4E79}
\definecolor{lightblue}{HTML}{EAF2F8}
\definecolor{darkgray}{HTML}{333333}
\definecolor{definitiongreen}{HTML}{EEF8EE}
\definecolor{definitionborder}{HTML}{8FB58F}
\definecolor{lightpink}{HTML}{FCE8EF}
\definecolor{warningyellow}{HTML}{FFF7D6}
\definecolor{warningborder}{HTML}{D5A500}

\setlength{\parindent}{0pt}
\setlength{\parskip}{0.65em}
\setlist[itemize]{topsep=0.25em,itemsep=0.15em}

% Dotted leaders in the table of contents.
\makeatletter
\def\@tocline#1#2#3#4#5#6#7{\relax
  \ifnum #1>\c@tocdepth
  \else
    \par\addpenalty\@secpenalty\addvspace{#2}%
    \begingroup\hyphenpenalty\@M
    \@ifempty{#4}{%
      \@tempdima\csname r@tocindent\number#1\endcsname\relax
    }{%
      \@tempdima#4\relax
    }%
    \parindent\z@\leftskip#3\relax\advance\leftskip\@tempdima\relax
    \rightskip\@pnumwidth plus4em\parfillskip-\@pnumwidth
    #5\leavevmode\hskip-\@tempdima
    \ifcase #1
      \or\or\hskip 1em\or\hskip 2em\else\hskip 3em
    \fi
    #6\nobreak\relax
    \dotfill\hbox to\@pnumwidth{\@tocpagenum{#7}}\par
    \nobreak
    \endgroup
  \fi}
\makeatother

\theoremstyle{definition}
\newtheorem{definition}{Definition}[section]
\tcolorboxenvironment{definition}{
  enhanced,
  breakable,
  colback=definitiongreen,
  colframe=definitionborder,
  boxrule=0.5pt,
  arc=1mm,
  left=2mm,
  right=2mm,
  top=1mm,
  bottom=1mm,
  before skip=0.8em,
  after skip=0.8em
}
\newtheorem{example}{Example}[section]
\theoremstyle{plain}
\newtheorem{theorem}{Theorem}[section]
\newtheorem{lemma}{Lemma}[section]
\newtheorem{proposition}{Proposition}[section]
\theoremstyle{remark}
\newtheorem*{remark}{Remark}
\newtheorem{exercise}[theorem]{Exercise}

\tcolorboxenvironment{exercise}{
	enhanced,
	colback=lightpink,
	colframe=WildStrawberry,
	boxrule=0.7pt,
	arc=2pt,
	left=6pt,
	right=6pt,
	top=5pt,
	bottom=5pt
}

\newtcolorbox{warningbox}{
  enhanced,
  breakable,
  colback=warningyellow,
  colframe=warningborder,
  coltitle=darkgray,
  fonttitle=\bfseries,
  title={Warning},
  boxrule=0.7pt,
  arc=2pt,
  left=6pt,
  right=6pt,
  top=5pt,
  bottom=5pt,
  before skip=0.8em,
  after skip=0.8em
}


\numberwithin{equation}{section}

\newcommand{\Z}{\mathbb{Z}}
\newcommand{\R}{\mathbb{R}}
\newcommand{\GL}{\operatorname{GL}}
\newcommand{\id}{\operatorname{id}}
\newcommand{\be}{\begin{eqnarray}}
\newcommand{\ee}{\end{eqnarray}}

\begin{document}

\begin{center}
  \colorbox{lightblue}{%
    \begin{minipage}{0.94\textwidth}
      \vspace{0.55em}
      \begin{tabular*}{\textwidth}{@{\extracolsep{\fill}}lr}
        {\color{courseblue}\bfseries\LARGE Lecture 3} &
        {\color{courseblue}\bfseries\LARGE MATH 327}\\[0.35em]
         & \textbf{Date: Sep. 9, 2026}\\
        & \textbf{Instructor: Manish Patnaik}
      \end{tabular*}
      \vspace{0.45em}
    \end{minipage}}
\end{center}

\setcounter{tocdepth}{1}
\tableofcontents

\section{Recollections on Permutation Groups}
 
 For a positive integer $n$, let us define
\be
[n]:=\{1,2,\ldots,n\}.
\ee
A permutation of $[n]$ is a bijection $p\colon[n]\to[n]$. It may be displayed
in two-line notation as
\be
p=
\begin{pmatrix}
1&2&\cdots&n\\
p(1)&p(2)&\cdots&p(n)
\end{pmatrix}.
\ee

If $p$ and $q$ are permutations, we use the convention
\be
qp=q\circ p;
\ee
thus the rightmost permutation is applied first. Let us also note that if $p$
is a permutation, we write $p^{-1}$ for the inverse permutation; that is, if
$p(i)=k$, then $p^{-1}(k)=i$. One can check that
\be
pp^{-1}=p^{-1}p=\id,
\ee
where $\id$ is the identity permutation.

Let $S_n$ denote the set of all permutations of $[n]$.

\begin{proposition}
	The triple $(S_n,\circ,\id)$ is a group.
\end{proposition}

\begin{proof}
	The composite of two bijections is again a bijection, so $S_n$ is closed
	under composition. Composition of functions is associative, the identity
	map $\id$ is the identity element, and the inverse of a bijection is again a
	bijection. Hence all the group axioms hold.
\end{proof}


\newpage
\section{Cycle Notation}
For future applications, we need to introduce the important `cycle notation'. 

\begin{definition}
Fix a positive integer $n$ and an integer $k$ with $1\leq k\leq n$.
For distinct elements $a_1,\ldots,a_k\in[n]$, the \emph{$k$-cycle}
$\sigma=(a_1\ \cdots\ a_k)$ is the permutation defined by
\be
\sigma(a_i)=a_{i+1}\quad(1\leq i<k),\qquad \sigma(a_k)=a_1.
\ee
Thus it sends $a_1\mapsto a_2\mapsto\cdots\mapsto a_k\mapsto a_1$
and fixes every element outside $\{a_1,\ldots,a_k\}$.
Every one-cycle represents the identity permutation, which we denote by
$\id$ or $(1)$. The empty product also represents the identity.
A cycle of the form $(a\ b)$, with $a\neq b$, is called a \emph{transposition}.
\end{definition}


For example, the cycle
\be
p=(1\ 2\ 3)
\ee
sends $1$ to $2$, $2$ to $3$, and $3$ to $1$, while fixing every other
element. Thus, as a permutation of $[5]$,
\be
(1\ 2\ 3)=
\begin{pmatrix}
1&2&3&4&5\\
2&3&1&4&5
\end{pmatrix}.
\ee

	Cyclic rotations of the entries represent the same cycle. For example,
	\be
	(1\ 2\ 3)=(2\ 3\ 1)=(3\ 1\ 2).
	\ee

\subsection*{Composing cycles} 

If $p, q$ are cycles then we would like to develop some tools for composing them, i.e. for calculating the product $pq$.  Let's start with a simple example: 

\begin{example}[Composing disjoint cycles]
Let $q=(2\ 3\ 4)$ and $p=(1\ 5)$. Then
\be
qp=(2\ 3\ 4)(1\ 5)
=
\begin{pmatrix}
1&2&3&4&5\\
5&3&4&2&1
\end{pmatrix}.
\ee
When the entries in cycles are disjoint, the two permutations commute:
\be
(2\ 3\ 4)(1\ 5)=(1\ 5)(2\ 3\ 4).
\ee

Cycles that are not disjoint \textit{need not commute.} For example,
you can verify that, as permutations of $[4]$,
$(2\ 4)(1\ 4)\neq(1\ 4)(2\ 4)$. Indeed, we have
\be
(2\ 4)(1\ 4) = \begin{pmatrix}
	1&2&3&4\\
	2&4&3&1
\end{pmatrix}, \qquad
(1\ 4)(2\ 4) = \begin{pmatrix}
	1&2&3&4\\
	4&1&3&2
\end{pmatrix}.
\ee

\end{example}


Here is a basic fact:

\begin{lemma} 
Every permutation is a product of disjoint cycles. \end{lemma}

\begin{proof}
Let $p$ be a permutation of $[n]$. Choose $a_1\in[n]$ and consider the sequence obtained by repeatedly
applying $p$ to $a_1$, namely,
\be
a_1,\ p(a_1),\ p^2(a_1),\ldots.
\ee
Since $[n]$ is finite, some term must eventually repeat. Let $k$ be the first
positive integer for which $p^k(a_1)$ is equal to an earlier term, say
$p^k(a_1)=p^j(a_1)$ with $0\leq j<k$. If $j>0$, applying $p^{-1}$ gives
$p^{k-1}(a_1)=p^{j-1}(a_1)$, contradicting the choice of $k$. Hence $j=0$.

It follows that
\be
a_1,\ p(a_1),\ldots,p^{k-1}(a_1)
\ee
are distinct and that $p$ acts on them as the cycle
\be
(a_1\ p(a_1)\ \cdots\ p^{k-1}(a_1)).
\ee

If this cycle does not contain every element of $[n]$, choose an unused
element $a_2$ and repeat the argument. The new cycle is disjoint from the
first one: if the two cycles shared an element, say $d$, then for some $i, j \geq 0$ we must have  \be d=p^i(a_1) = p^j(a_2). \ee
Repeatedly applying $p$ (or $p^{-1}$) to this equation shows that
$a_2$ lies in the first cycle, a contradiction.

Repeating this process gives pairwise disjoint cycles containing
every element of $[n]$. Their product agrees with $p$ on every element, and
therefore equals $p$. (Cycles of length one may be omitted.)

\end{proof}

\begin{example}[Disjoint cycles in $S_5$]
Consider
\be
p=
\begin{pmatrix}
1&2&3&4&5\\
4&1&5&2&3
\end{pmatrix}.
\ee
Starting with $1$ and following its images gives
\be
1\longmapsto4\longmapsto2\longmapsto1,
\ee
while the remaining elements that are not yet accounted for satisfy
\be
3\longmapsto5\longmapsto3.
\ee
Thus the decomposition of $p$ into disjoint cycles is
\be
p=(1\ 4\ 2)(3\ 5).
\ee

\end{example}



\end{document}
